\section{Preliminaries and Problem Setting}
\label{sec:prelim}

\paragraph{Cross-silo federated graph learning.}
There are $K$ institutions (clients). Client $k$ holds a private graph
$\mathcal{G}_k=(\mathcal{V}_k,\mathcal{E}_k,\mathbf{X}_k,\mathbf{y}_k,\mathbf{s}_k)$
with node features $\mathbf{X}_k\in\R^{n_k\times d}$, binary task labels
$\mathbf{y}_k\in\{0,1\}^{n_k}$ (e.g.\ fraud/default/recidivism), and a binary
\emph{sensitive attribute} $\mathbf{s}_k\in\{0,1\}^{n_k}$ (e.g.\ gender, age,
race). Raw graphs never leave the client. A server coordinates $T$
communication rounds; each round clients perform $E$ local epochs and transmit
model parameters, which the server aggregates. We evaluate the global model on
the pooled held-out test nodes of all clients. Statistical heterogeneity is
modelled by a symmetric Dirichlet($\alpha$) split over labels, the standard FL
non-IID model \cite{hsu2019measuring}; smaller $\alpha$ is more skewed.

\paragraph{Group fairness.}
For predicted probabilities $\hat{y}$ and sensitive attribute $s$ we report
three group-fairness metrics (lower is fairer):
\begin{align}
\dpd &= \big|\,\E[\hat{y}\mid s{=}0]-\E[\hat{y}\mid s{=}1]\,\big|
        &&\text{(Demographic Parity Difference)},\\
\eod &= \big|\,\mathrm{TPR}_{s=0}-\mathrm{TPR}_{s=1}\,\big|
        &&\text{(Equal Opportunity Difference)},\\
\Delta_{\mathrm{EOdds}} &= \max\big(|\Delta\mathrm{TPR}|,\,|\Delta\mathrm{FPR}|\big)
        &&\text{(Equalized Odds)} .
\end{align}
$\dpd$ is computed on soft scores (matching the differentiable training term);
$\eod$ and Equalized Odds use thresholded predictions.

\paragraph{Differential privacy.}
A randomised mechanism $\mathcal{M}$ is $(\epsilon,\delta)$-differentially
private if for all adjacent datasets $D,D'$ and all outputs $S$,
$\Pr[\mathcal{M}(D)\in S]\le e^{\epsilon}\Pr[\mathcal{M}(D')\in S]+\delta$
\cite{dwork2014algorithmic}. We protect the \emph{sensitive attribute}:
adjacency is defined by flipping one node's $s_i$. The Gaussian mechanism adds
$\mathcal{N}(0,\sigma^2)$ noise calibrated to the $L_2$ sensitivity of the
released quantity; composition across rounds is tracked with R\'enyi DP
\cite{mironov2017renyi}. For the Gaussian mechanism with noise multiplier
$z=\sigma/\Delta$, the R\'enyi divergence of order $\alpha$ is
$\epsilon_{\mathrm{RDP}}(\alpha)=\alpha/(2z^2)$; over $T$ releases it composes as
$T\alpha/(2z^2)$, and the tight conversion to $(\epsilon,\delta)$-DP is
$\epsilon=\min_{\alpha>1}\big[\epsilon_{\mathrm{RDP}}(\alpha)+\log(1/\delta)/(\alpha-1)\big]$.

\paragraph{Threat model for robustness.}
Up to $f<K/2$ clients are Byzantine and may send arbitrary updates. We consider
data poisoning (label flipping), classical model poisoning (Gaussian, sign-flip,
scaling/model-replacement \cite{bagdasaryan2020backdoor}), the strong
omniscient attacks ALIE \cite{baruch2019little} and IPM \cite{xie2019fall},
and a \emph{fairness-poisoning} attack: a malicious client transmits a biased,
scaled update while reporting fabricated fairness/utility statistics
($\dpd\!\approx\!0$, high AUC) to capture a fairness-aware aggregator. The
server is honest but curious and never accesses raw data.

\paragraph{Objective.}
We seek a global model that jointly (i) maximises task utility (AUC-ROC), (ii)
minimises group unfairness, (iii) satisfies a target $(\epsilon,\delta)$-DP for
the sensitive attribute, and (iv) retains utility and fairness under Byzantine
corruption---the trustworthiness trade-off that motivates
Section~\ref{sec:method}.
